% CAS BMAI HS26, weekend 1: the cold open, Szymborska variant.
% "Guess, Check, Correct".
%
% WHAT THIS DECK IS. It is the Popper cold open, ../../cold-open-weekend1-popper/,
% with ONE FRAME ADDED AT THE FRONT: Wislawa Szymborska's 1996 Nobel lecture.
% Frames 2 to 6 are that deck byte for byte apart from their banner numbers and
% three added lines on the references frame. Built 2026-09-01 on Carlos's
% request, after a search for a literary cold open turned up four candidates
% (Paz, Tagore, Szymborska, Beckett) and he chose Szymborska.
% DIFFING AGAINST THE POPPER DECK IS THEREFORE MEANINGFUL. Keep it that way:
% if you change a shared frame here, change it there in the same pass, or say
% in this header why the two have been allowed to diverge.
%
% Light projection-first Metropolis with the Libertinus serif, inherited from
% beamer-preamble.tex in this directory (copied on 2026-08-31 from the sibling
% Descartes deck ../../cold-open-weekend1/slides/, which had already been diffed
% against the lecture-slides skill; the diff is recorded at the head of that file).
% NOTE, and it is an open decision rather than an oversight: the HS26 house style
% was changed to Paper and Grid on 2026-09-01 and the 08:00, 09:00 and 11:00
% decks were restyled to it. All three cold opens, this one included, are still
% on this earlier light Metropolis. See ../README.md, "Open decisions".
%
% Engine: pdflatex, TWICE.
% Build:  cd slides && pdflatex slides.tex && pdflatex slides.tex
%
% STAGING. Six frames. Frames 1 to 5 are the 08:00 opening on Friday
% 4 September 2026, inside the "Intro and onboarding" block, before anything
% administrative. About three and a half minutes. Frame 6 is a reference frame:
% it is a handout page and a rights record, and it is NOT projected during those
% minutes. Cut it from the projected run or leave it as the page the deck
% happens to end on, but do not read it aloud.
%
% THE ARGUMENT. Every algorithm on Friday morning is one loop: guess, check
% against something you can actually compute, read the signed miss, correct,
% repeat. Szymborska describes that loop from the inside, as the thing a poet
% does to a poem, and she names the state it starts from: "I don't know".
% Popper then answers her in the same words, "We do not know", and adds the
% guess. Popper wrote the loop down as a formula in 1972 and defended it as the
% whole method of science. Turing said the same thing about machines that learn,
% one year before Robbins and Monro. Whewell said it in 1858 about inventive
% minds in general, and threw in overshooting for free. One poet, three
% witnesses, one loop, and the room has agreed with all four before a single
% line of code is written.
%
% WHY A POET FIRST. The three witnesses all make a claim ABOUT method, from
% outside it. Szymborska makes no claim at all. She reports what it feels like
% to finish something, read it back, and find it will not hold, which is an
% experience every person in the room has had and none of them calls a method.
% The frame costs about forty seconds and it buys the other three an audience
% that is already nodding.
%
% WHY POPPER RATHER THAN DESCARTES. The sibling deck opens on Descartes, and
% its own design spec records the objection: Descartes wants certainty by
% deduction from ideas so clear that doubt is impossible, and weekend 1 offers
% a guess, a step that can overshoot, and a bound rather than a proof. The
% Descartes deck answers that by saying the inversion out loud. Popper needs no
% such repair, because Popper's method IS this weekend's method.
%
% HOUSE RULES OBEYED. Dry factual frame titles. Original language above the
% gloss (see THE ORIGINAL-LANGUAGE PROBLEM below). Source line in muted small
% type on every quotation frame. One plate per thinker frame, framed and
% captioned in one line. NO CLOSING PUNCHLINE and no \takeaway. ASCII
% punctuation only: there is no prose "---" and no prose "--" anywhere in this
% file, and the TikZ chain is drawn with "to" rather than "--" so that even the
% diagram is dash free.
% See /Users/carloscotrini/Documents/git_fdd/projects/.claude/skills/lecture-slides/SKILL.md
%
% THE ORIGINAL-LANGUAGE PROBLEM, and how this deck handles it.
% Carlos asked for the original language above the translation on every frame.
% Three of the five quotation frames carry text WRITTEN IN ENGLISH: Popper's
% Conjectures and Refutations (1963) and Objective Knowledge (1972), Turing's
% Mind paper (1950), and Whewell (1858). For those the original IS the English
% on the slide, and the muted line underneath is therefore a PARAPHRASE rather
% than a translation. Each such source line says so in as many words, so that
% nobody mistakes a gloss for a translation.
% Frame 1 is the one frame where the rule works exactly as written: Szymborska
% delivered the lecture in Polish, nobelprize.org carries the Polish and the
% English side by side, both were pinned, so the Polish is the body text and the
% Baranczak and Cavanagh English sits under it as a real translation.
% The one genuinely German-original quotation, "We do not know: we can only
% guess", could NOT be pinned in the build session: two different German
% wordings circulate and neither was traceable to a printed page of Logik der
% Forschung. Its original-language slot on frame 2 therefore draws a dashed
% marker naming the work and the locus, via \origslot in the preamble. That is
% the same principle as a missing plate. Do not fill it in from a quotation
% site.
%
% PROVENANCE, and it is NOT uniform across the quotations.
%   PINNED verbatim in the research session against complete texts:
%     both Szymborska passages, Polish AND English (frame 1),
%     the Conjectures and Refutations preface sentence (frame 3),
%     the whole tetradic-schema passage behind frame 3's diagram,
%     both Turing sentences (frame 4),
%     both Whewell sentences (frame 5).
%   LOCATED but NOT pinned, and carrying a known punctuation variant:
%     "We do not know: we can only guess" (frame 2). Colon and semicolon
%     versions both circulate.
% The full table is in ../SOURCES.md. Read it before projecting anything.
%
% IMAGES. Three of the four plates this deck wants are PRESENT in img/ and were
% licence-checked at the source and inspected: turing-1951.jpg (public domain),
% whewell-wellcome.jpg (CC BY 4.0) and szymborska-kubik.jpg (CC BY 3.0, added
% 2026-09-01). img/IMAGE-SOURCES.txt records artist, date, licence and Commons
% URL for each, and \plateh draws a dashed box naming any file that is missing
% so the gap is visible rather than silent.
% ONE PLATE IS STILL ABSENT AND ITS BOX STILL PROJECTS: img/popper-lse.jpg. It
% carries a rights question that only Carlos can settle, because Popper is
% protected in Switzerland until the end of 2064 and the Commons record offers
% only an institutional "no known copyright restrictions" statement rather than
% a licence. Read that file before filling the slot.
% Do not "fix" a dashed box by committing a stand-in image: a placeholder that
% looks like a picture is worse than a box saying a picture is missing.

\documentclass[aspectratio=169]{beamer}
% ===== CAS BMAI HS26, weekend 1 cold open (Popper variant): Beamer preamble =====
% Light, projection-first Metropolis with the Libertinus serif.
%
% PROVENANCE. Copied on 2026-08-31 from the sibling deck
%   ../../cold-open-weekend1/slides/beamer-preamble.tex
% which was itself copied from
%   /Users/carloscotrini/Documents/git_fdd/projects/.claude/skills/lecture-slides/templates/beamer-preamble.tex
% on 2026-08-26, then diffed against
%   /Users/carloscotrini/Documents/git_fdd/projects/.claude/skills/lecture-slides/SKILL.md
% before a single frame was written, which is what that skill instructs
% ("When you inherit a preamble, diff its section page and title page templates
% against what this skill prescribes, and strip anything extra before writing a
% single frame").
%
% WHAT THE DIFF FOUND, and what was done about it:
%   1. \takeaway            ABSENT from the template already. Still absent here,
%                           deliberately. See the closing comment block.
%   2. \sectionsub, \cursub ABSENT from the template already. Still absent.
%                           These are the "running question" mechanism that the
%                           skill records creeping back into decks by wholesale
%                           preamble copying. Do not add them.
%   3. section page,
%      title page templates The template defines NEITHER, because plain Metropolis
%                           supplies both and this deck calls neither: it has no
%                           \section and no \maketitle. Nothing to strip.
%   4. good / bad colours   The template already ships Okabe-Ito blue #0072B2 and
%                           vermillion #D55E00. Kept verbatim. No red/green pair
%                           exists in this file, and this deck uses neither colour.
%   5. muted                Kept at black!55, the legible floor. Do not lighten it:
%                           the muted row of the mapping table is the whole point
%                           of this deck and it has to survive a projector.
%
% ADDED FOR THIS DECK, and nothing else (each addition is marked ADDED below):
%   \plateh    a framed plate that degrades to a visible placeholder box naming
%              the file, so the deck compiles with an empty img/ directory.
%   \framedimgh  the hairline image frame the plate uses, lifted from
%              /Users/carloscotrini/Documents/git_fdd/projects/mandatory-weekends-opening/slides/beamer-preamble.tex
%   \origslot  ADDED 2026-08-31 FOR THIS DECK ONLY. The original-language slot
%              of a quotation frame, which degrades to a visible dashed marker
%              when the original has not been pinned yet. Same principle as
%              \plateh: a gap that is legible on the slide cannot be shipped by
%              accident. See the definition at the foot of this file.
%
% Engine: pdflatex, twice. No fontspec, no lualatex. The cold open quotes French
% only, so nothing here needs a Unicode engine. (A Greek candidate would have:
% see COLD-OPEN-CANDIDATES.md section 0.)
%
% Each deck does:
%   \documentclass[aspectratio=169]{beamer}
%   % ===== CAS BMAI HS26, weekend 1 cold open (Popper variant): Beamer preamble =====
% Light, projection-first Metropolis with the Libertinus serif.
%
% PROVENANCE. Copied on 2026-08-31 from the sibling deck
%   ../../cold-open-weekend1/slides/beamer-preamble.tex
% which was itself copied from
%   /Users/carloscotrini/Documents/git_fdd/projects/.claude/skills/lecture-slides/templates/beamer-preamble.tex
% on 2026-08-26, then diffed against
%   /Users/carloscotrini/Documents/git_fdd/projects/.claude/skills/lecture-slides/SKILL.md
% before a single frame was written, which is what that skill instructs
% ("When you inherit a preamble, diff its section page and title page templates
% against what this skill prescribes, and strip anything extra before writing a
% single frame").
%
% WHAT THE DIFF FOUND, and what was done about it:
%   1. \takeaway            ABSENT from the template already. Still absent here,
%                           deliberately. See the closing comment block.
%   2. \sectionsub, \cursub ABSENT from the template already. Still absent.
%                           These are the "running question" mechanism that the
%                           skill records creeping back into decks by wholesale
%                           preamble copying. Do not add them.
%   3. section page,
%      title page templates The template defines NEITHER, because plain Metropolis
%                           supplies both and this deck calls neither: it has no
%                           \section and no \maketitle. Nothing to strip.
%   4. good / bad colours   The template already ships Okabe-Ito blue #0072B2 and
%                           vermillion #D55E00. Kept verbatim. No red/green pair
%                           exists in this file, and this deck uses neither colour.
%   5. muted                Kept at black!55, the legible floor. Do not lighten it:
%                           the muted row of the mapping table is the whole point
%                           of this deck and it has to survive a projector.
%
% ADDED FOR THIS DECK, and nothing else (each addition is marked ADDED below):
%   \plateh    a framed plate that degrades to a visible placeholder box naming
%              the file, so the deck compiles with an empty img/ directory.
%   \framedimgh  the hairline image frame the plate uses, lifted from
%              /Users/carloscotrini/Documents/git_fdd/projects/mandatory-weekends-opening/slides/beamer-preamble.tex
%   \origslot  ADDED 2026-08-31 FOR THIS DECK ONLY. The original-language slot
%              of a quotation frame, which degrades to a visible dashed marker
%              when the original has not been pinned yet. Same principle as
%              \plateh: a gap that is legible on the slide cannot be shipped by
%              accident. See the definition at the foot of this file.
%
% Engine: pdflatex, twice. No fontspec, no lualatex. The cold open quotes French
% only, so nothing here needs a Unicode engine. (A Greek candidate would have:
% see COLD-OPEN-CANDIDATES.md section 0.)
%
% Each deck does:
%   \documentclass[aspectratio=169]{beamer}
%   % ===== CAS BMAI HS26, weekend 1 cold open (Popper variant): Beamer preamble =====
% Light, projection-first Metropolis with the Libertinus serif.
%
% PROVENANCE. Copied on 2026-08-31 from the sibling deck
%   ../../cold-open-weekend1/slides/beamer-preamble.tex
% which was itself copied from
%   /Users/carloscotrini/Documents/git_fdd/projects/.claude/skills/lecture-slides/templates/beamer-preamble.tex
% on 2026-08-26, then diffed against
%   /Users/carloscotrini/Documents/git_fdd/projects/.claude/skills/lecture-slides/SKILL.md
% before a single frame was written, which is what that skill instructs
% ("When you inherit a preamble, diff its section page and title page templates
% against what this skill prescribes, and strip anything extra before writing a
% single frame").
%
% WHAT THE DIFF FOUND, and what was done about it:
%   1. \takeaway            ABSENT from the template already. Still absent here,
%                           deliberately. See the closing comment block.
%   2. \sectionsub, \cursub ABSENT from the template already. Still absent.
%                           These are the "running question" mechanism that the
%                           skill records creeping back into decks by wholesale
%                           preamble copying. Do not add them.
%   3. section page,
%      title page templates The template defines NEITHER, because plain Metropolis
%                           supplies both and this deck calls neither: it has no
%                           \section and no \maketitle. Nothing to strip.
%   4. good / bad colours   The template already ships Okabe-Ito blue #0072B2 and
%                           vermillion #D55E00. Kept verbatim. No red/green pair
%                           exists in this file, and this deck uses neither colour.
%   5. muted                Kept at black!55, the legible floor. Do not lighten it:
%                           the muted row of the mapping table is the whole point
%                           of this deck and it has to survive a projector.
%
% ADDED FOR THIS DECK, and nothing else (each addition is marked ADDED below):
%   \plateh    a framed plate that degrades to a visible placeholder box naming
%              the file, so the deck compiles with an empty img/ directory.
%   \framedimgh  the hairline image frame the plate uses, lifted from
%              /Users/carloscotrini/Documents/git_fdd/projects/mandatory-weekends-opening/slides/beamer-preamble.tex
%   \origslot  ADDED 2026-08-31 FOR THIS DECK ONLY. The original-language slot
%              of a quotation frame, which degrades to a visible dashed marker
%              when the original has not been pinned yet. Same principle as
%              \plateh: a gap that is legible on the slide cannot be shipped by
%              accident. See the definition at the foot of this file.
%
% Engine: pdflatex, twice. No fontspec, no lualatex. The cold open quotes French
% only, so nothing here needs a Unicode engine. (A Greek candidate would have:
% see COLD-OPEN-CANDIDATES.md section 0.)
%
% Each deck does:
%   \documentclass[aspectratio=169]{beamer}
%   \input{beamer-preamble}
% and is compiled FROM the directory that holds this file so \input resolves.
%
% House rule: ASCII punctuation only. No prose "---" (renders an em dash) and
% no prose "--" (renders an en dash). The ONLY allowed "--" is a TikZ line
% segment, e.g. \draw (a) -- (b). Use a comma, colon, parentheses, or "to".

\usepackage[T1]{fontenc}
\usepackage{libertinus}   % Libertinus Serif + Biolinum + matching Libertinus Math
\usefonttheme{serif}      % body, headings, and math all in the serif

\usetheme{metropolis}
\metroset{
  sectionpage=progressbar,
  numbering=fraction,
  progressbar=foot,
  block=fill
}

% --- accents for a LIGHT canvas (same names; retune hues per course) ---
\definecolor{accent}{HTML}{0369A1}    % deep sky blue (primary accent)
\definecolor{accent2}{HTML}{B45309}   % dark amber (secondary highlight)
\definecolor{good}{HTML}{0072B2}   % Okabe-Ito blue: a win, the chosen option
\definecolor{bad}{HTML}{D55E00}    % Okabe-Ito vermillion. NO red/green: see SKILL.md
\colorlet{panel}{black!8}             % a panel slightly darker than the bg

% MANDATORY: white background, near-black text (the projection-first default).
\setbeamercolor{normal text}{fg=black!87, bg=white}
\setbeamercolor{alerted text}{fg=accent}
\setbeamercolor{progress bar in head/foot}{fg=accent}
\setbeamercolor{progress bar in section page}{fg=accent}
\setbeamercolor{title separator}{fg=accent}
% MANDATORY: frametitle as a pale panel strip with near-black text.
\setbeamercolor{frametitle}{fg=black!88, bg=panel}
\setbeamercolor{frametitle right}{fg=black!88, bg=panel}

% --- math: Libertinus already supplies it; amssymb/dsfont for extras ---
\usepackage{amsmath, amssymb}
\usepackage{dsfont}   % \mathds{1} indicator (amssymb has no blackboard-bold digit 1)

% --- TikZ, styled dark-on-light so every diagram reads on a projector ---
\usepackage{tikz}
% overlay-beamer-styles provides "visible on=<n->", which the no-shift rule
% (SKILL.md Pattern F) requires instead of wrapping a node in \only. Without it
% every deck built from this template fails with "I do not know the key
% '/tikz/visible on'".
\usetikzlibrary{arrows.meta, positioning, calc, backgrounds, fit, shapes.geometric,
                overlay-beamer-styles}
\tikzset{>={Stealth[round]}}
\colorlet{fg}{black!85}
% MANDATORY: keep muted DARK enough on the light background. black!55 is the
% legible floor; anything lighter washes out under a projector. Do not raise it.
\colorlet{muted}{black!55}
\tikzset{
  every node/.style={text=fg, font=\footnotesize},
  edge/.style={draw=fg, thick, ->},
  dimedge/.style={draw=muted, thick, ->},
  nodebox/.style={draw=fg, thick, rounded corners=2pt, inner sep=5pt, text=fg},
  accentbox/.style={draw=accent, line width=1pt, rounded corners=2pt, inner sep=5pt, text=fg},
  gridcell/.style={draw=muted, minimum size=1.05cm, anchor=center, text=fg},
  treenode/.style={circle, draw=fg, thick, minimum size=0.9cm, inner sep=1pt, text=fg},
  good fill/.style={fill=good!14!white, draw=good, text=fg},
  bad fill/.style={fill=bad!12!white, draw=bad, text=fg},
  acc fill/.style={fill=accent!13!white, draw=accent, text=fg},
  % ADDED: the dashed slot a missing plate leaves behind.
  plateslot/.style={draw=muted, dashed, line width=0.8pt, rounded corners=4pt,
                    align=center, text=muted, font=\scriptsize, inner sep=6pt},
}

% ---------------------------------------------------------------------------
% ADDED for this deck: framed plate, height constrained.
% Lifted from mandatory-weekends-opening/slides/beamer-preamble.tex, where it is
% called \framedimgh and used for the Kant and Friedrich plates.
% ---------------------------------------------------------------------------
\newcommand{\framedimgh}[2]{%
  \tikz{\node[inner sep=0pt, draw=black!18, line width=0.8pt, rounded corners=4pt]%
        {\includegraphics[height=#1]{#2}};}}

% ---------------------------------------------------------------------------
% ADDED for this deck: \plateh{<height>}{<placeholder width>}{<path>}
%
% If the file is present, this is exactly \framedimgh. If it is ABSENT, it draws
% a dashed box of the same height naming the file that belongs there, so the deck
% compiles and projects with an empty img/ directory and the gap is visible
% rather than silent.
%
% This exists because no image is downloaded into this repository: the licences
% are recorded in img/IMAGE-SOURCES.txt and the files are fetched by hand, once,
% by whoever is going to project the deck. Do not "fix" this by committing a
% stand-in image; a placeholder that looks like a picture is worse than a box
% that says a picture is missing.
% ---------------------------------------------------------------------------
\newcommand{\plateh}[3]{%
  \IfFileExists{#3}{\framedimgh{#1}{#3}}{%
    \tikz{\node[plateslot, minimum height=#1, minimum width=#2]%
          {Plate not present\\[2pt]\texttt{#3}\\[2pt] See \texttt{img/IMAGE-SOURCES.txt}};}}}

% ---------------------------------------------------------------------------
% ADDED for this deck: \origslot{<original-language text>}{<work and locus>}
%
% House grammar rule 1 puts the ORIGINAL LANGUAGE on top and the English gloss
% underneath. House rule "do not put an unpinned quotation on a slide" says the
% original may only be projected once somebody has pinned it character for
% character against a primary text.
%
% Those two rules collide on this deck, because its one German-original
% quotation could NOT be pinned in the build session: two different German
% wordings circulate and neither was traceable to a printed page of
% Logik der Forschung. So the original-language slot takes the pinned text when
% there is one, and otherwise draws a dashed marker naming the work and the
% locus that still has to be checked.
%
% Pass a non-empty first argument and it renders as an ordinary quotation.
% Pass an empty first argument and the gap is visible on the slide.
%
% Do NOT "fix" an empty slot by pasting in a wording found in a quotation site.
% That is the exact failure the marker exists to prevent.
% ---------------------------------------------------------------------------
\newcommand{\origslot}[2]{%
  \if\relax\detokenize{#1}\relax
    \tikz{\node[plateslot, align=left, text width=0.86\textwidth]%
          {Original not yet pinned: #2.\\
           Pin it against a printed page before projecting, then pass it to
           \texttt{\string\origslot}. See \texttt{../SOURCES.md}.};}%
  \else
    #1%
  \fi}

% --- No closing-punchline helper here, on purpose: a frame ends on its final
%     reveal. Do not add one; a macro in the preamble reads as an instruction.
%
%     DELIBERATELY NOT DEFINED: \takeaway. The Weekend 1 FDD preamble that the
%     mandatory-weekends opening was copied from defines a centred amber-ruled
%     pull quote, and it predates the rule. Carlos, 2026-08-11: the big centred
%     blue punchline is out. A frame ends on its final reveal, and a closing
%     thought that is real content is the last ordinary line, in body text.
%     See .claude/skills/lecture-slides/SKILL.md, "No closing punchline".
%     The two closing thoughts in this deck (the placement line on frame 1 and
%     the "this hour is about the fourth" line on frame 4) are \footnotesize
%     body text at the foot of their frame, in ink, never in accent. ---

% and is compiled FROM the directory that holds this file so \input resolves.
%
% House rule: ASCII punctuation only. No prose "---" (renders an em dash) and
% no prose "--" (renders an en dash). The ONLY allowed "--" is a TikZ line
% segment, e.g. \draw (a) -- (b). Use a comma, colon, parentheses, or "to".

\usepackage[T1]{fontenc}
\usepackage{libertinus}   % Libertinus Serif + Biolinum + matching Libertinus Math
\usefonttheme{serif}      % body, headings, and math all in the serif

\usetheme{metropolis}
\metroset{
  sectionpage=progressbar,
  numbering=fraction,
  progressbar=foot,
  block=fill
}

% --- accents for a LIGHT canvas (same names; retune hues per course) ---
\definecolor{accent}{HTML}{0369A1}    % deep sky blue (primary accent)
\definecolor{accent2}{HTML}{B45309}   % dark amber (secondary highlight)
\definecolor{good}{HTML}{0072B2}   % Okabe-Ito blue: a win, the chosen option
\definecolor{bad}{HTML}{D55E00}    % Okabe-Ito vermillion. NO red/green: see SKILL.md
\colorlet{panel}{black!8}             % a panel slightly darker than the bg

% MANDATORY: white background, near-black text (the projection-first default).
\setbeamercolor{normal text}{fg=black!87, bg=white}
\setbeamercolor{alerted text}{fg=accent}
\setbeamercolor{progress bar in head/foot}{fg=accent}
\setbeamercolor{progress bar in section page}{fg=accent}
\setbeamercolor{title separator}{fg=accent}
% MANDATORY: frametitle as a pale panel strip with near-black text.
\setbeamercolor{frametitle}{fg=black!88, bg=panel}
\setbeamercolor{frametitle right}{fg=black!88, bg=panel}

% --- math: Libertinus already supplies it; amssymb/dsfont for extras ---
\usepackage{amsmath, amssymb}
\usepackage{dsfont}   % \mathds{1} indicator (amssymb has no blackboard-bold digit 1)

% --- TikZ, styled dark-on-light so every diagram reads on a projector ---
\usepackage{tikz}
% overlay-beamer-styles provides "visible on=<n->", which the no-shift rule
% (SKILL.md Pattern F) requires instead of wrapping a node in \only. Without it
% every deck built from this template fails with "I do not know the key
% '/tikz/visible on'".
\usetikzlibrary{arrows.meta, positioning, calc, backgrounds, fit, shapes.geometric,
                overlay-beamer-styles}
\tikzset{>={Stealth[round]}}
\colorlet{fg}{black!85}
% MANDATORY: keep muted DARK enough on the light background. black!55 is the
% legible floor; anything lighter washes out under a projector. Do not raise it.
\colorlet{muted}{black!55}
\tikzset{
  every node/.style={text=fg, font=\footnotesize},
  edge/.style={draw=fg, thick, ->},
  dimedge/.style={draw=muted, thick, ->},
  nodebox/.style={draw=fg, thick, rounded corners=2pt, inner sep=5pt, text=fg},
  accentbox/.style={draw=accent, line width=1pt, rounded corners=2pt, inner sep=5pt, text=fg},
  gridcell/.style={draw=muted, minimum size=1.05cm, anchor=center, text=fg},
  treenode/.style={circle, draw=fg, thick, minimum size=0.9cm, inner sep=1pt, text=fg},
  good fill/.style={fill=good!14!white, draw=good, text=fg},
  bad fill/.style={fill=bad!12!white, draw=bad, text=fg},
  acc fill/.style={fill=accent!13!white, draw=accent, text=fg},
  % ADDED: the dashed slot a missing plate leaves behind.
  plateslot/.style={draw=muted, dashed, line width=0.8pt, rounded corners=4pt,
                    align=center, text=muted, font=\scriptsize, inner sep=6pt},
}

% ---------------------------------------------------------------------------
% ADDED for this deck: framed plate, height constrained.
% Lifted from mandatory-weekends-opening/slides/beamer-preamble.tex, where it is
% called \framedimgh and used for the Kant and Friedrich plates.
% ---------------------------------------------------------------------------
\newcommand{\framedimgh}[2]{%
  \tikz{\node[inner sep=0pt, draw=black!18, line width=0.8pt, rounded corners=4pt]%
        {\includegraphics[height=#1]{#2}};}}

% ---------------------------------------------------------------------------
% ADDED for this deck: \plateh{<height>}{<placeholder width>}{<path>}
%
% If the file is present, this is exactly \framedimgh. If it is ABSENT, it draws
% a dashed box of the same height naming the file that belongs there, so the deck
% compiles and projects with an empty img/ directory and the gap is visible
% rather than silent.
%
% This exists because no image is downloaded into this repository: the licences
% are recorded in img/IMAGE-SOURCES.txt and the files are fetched by hand, once,
% by whoever is going to project the deck. Do not "fix" this by committing a
% stand-in image; a placeholder that looks like a picture is worse than a box
% that says a picture is missing.
% ---------------------------------------------------------------------------
\newcommand{\plateh}[3]{%
  \IfFileExists{#3}{\framedimgh{#1}{#3}}{%
    \tikz{\node[plateslot, minimum height=#1, minimum width=#2]%
          {Plate not present\\[2pt]\texttt{#3}\\[2pt] See \texttt{img/IMAGE-SOURCES.txt}};}}}

% ---------------------------------------------------------------------------
% ADDED for this deck: \origslot{<original-language text>}{<work and locus>}
%
% House grammar rule 1 puts the ORIGINAL LANGUAGE on top and the English gloss
% underneath. House rule "do not put an unpinned quotation on a slide" says the
% original may only be projected once somebody has pinned it character for
% character against a primary text.
%
% Those two rules collide on this deck, because its one German-original
% quotation could NOT be pinned in the build session: two different German
% wordings circulate and neither was traceable to a printed page of
% Logik der Forschung. So the original-language slot takes the pinned text when
% there is one, and otherwise draws a dashed marker naming the work and the
% locus that still has to be checked.
%
% Pass a non-empty first argument and it renders as an ordinary quotation.
% Pass an empty first argument and the gap is visible on the slide.
%
% Do NOT "fix" an empty slot by pasting in a wording found in a quotation site.
% That is the exact failure the marker exists to prevent.
% ---------------------------------------------------------------------------
\newcommand{\origslot}[2]{%
  \if\relax\detokenize{#1}\relax
    \tikz{\node[plateslot, align=left, text width=0.86\textwidth]%
          {Original not yet pinned: #2.\\
           Pin it against a printed page before projecting, then pass it to
           \texttt{\string\origslot}. See \texttt{../SOURCES.md}.};}%
  \else
    #1%
  \fi}

% --- No closing-punchline helper here, on purpose: a frame ends on its final
%     reveal. Do not add one; a macro in the preamble reads as an instruction.
%
%     DELIBERATELY NOT DEFINED: \takeaway. The Weekend 1 FDD preamble that the
%     mandatory-weekends opening was copied from defines a centred amber-ruled
%     pull quote, and it predates the rule. Carlos, 2026-08-11: the big centred
%     blue punchline is out. A frame ends on its final reveal, and a closing
%     thought that is real content is the last ordinary line, in body text.
%     See .claude/skills/lecture-slides/SKILL.md, "No closing punchline".
%     The two closing thoughts in this deck (the placement line on frame 1 and
%     the "this hour is about the fourth" line on frame 4) are \footnotesize
%     body text at the foot of their frame, in ink, never in accent. ---

% and is compiled FROM the directory that holds this file so \input resolves.
%
% House rule: ASCII punctuation only. No prose "---" (renders an em dash) and
% no prose "--" (renders an en dash). The ONLY allowed "--" is a TikZ line
% segment, e.g. \draw (a) -- (b). Use a comma, colon, parentheses, or "to".

\usepackage[T1]{fontenc}
\usepackage{libertinus}   % Libertinus Serif + Biolinum + matching Libertinus Math
\usefonttheme{serif}      % body, headings, and math all in the serif

\usetheme{metropolis}
\metroset{
  sectionpage=progressbar,
  numbering=fraction,
  progressbar=foot,
  block=fill
}

% --- accents for a LIGHT canvas (same names; retune hues per course) ---
\definecolor{accent}{HTML}{0369A1}    % deep sky blue (primary accent)
\definecolor{accent2}{HTML}{B45309}   % dark amber (secondary highlight)
\definecolor{good}{HTML}{0072B2}   % Okabe-Ito blue: a win, the chosen option
\definecolor{bad}{HTML}{D55E00}    % Okabe-Ito vermillion. NO red/green: see SKILL.md
\colorlet{panel}{black!8}             % a panel slightly darker than the bg

% MANDATORY: white background, near-black text (the projection-first default).
\setbeamercolor{normal text}{fg=black!87, bg=white}
\setbeamercolor{alerted text}{fg=accent}
\setbeamercolor{progress bar in head/foot}{fg=accent}
\setbeamercolor{progress bar in section page}{fg=accent}
\setbeamercolor{title separator}{fg=accent}
% MANDATORY: frametitle as a pale panel strip with near-black text.
\setbeamercolor{frametitle}{fg=black!88, bg=panel}
\setbeamercolor{frametitle right}{fg=black!88, bg=panel}

% --- math: Libertinus already supplies it; amssymb/dsfont for extras ---
\usepackage{amsmath, amssymb}
\usepackage{dsfont}   % \mathds{1} indicator (amssymb has no blackboard-bold digit 1)

% --- TikZ, styled dark-on-light so every diagram reads on a projector ---
\usepackage{tikz}
% overlay-beamer-styles provides "visible on=<n->", which the no-shift rule
% (SKILL.md Pattern F) requires instead of wrapping a node in \only. Without it
% every deck built from this template fails with "I do not know the key
% '/tikz/visible on'".
\usetikzlibrary{arrows.meta, positioning, calc, backgrounds, fit, shapes.geometric,
                overlay-beamer-styles}
\tikzset{>={Stealth[round]}}
\colorlet{fg}{black!85}
% MANDATORY: keep muted DARK enough on the light background. black!55 is the
% legible floor; anything lighter washes out under a projector. Do not raise it.
\colorlet{muted}{black!55}
\tikzset{
  every node/.style={text=fg, font=\footnotesize},
  edge/.style={draw=fg, thick, ->},
  dimedge/.style={draw=muted, thick, ->},
  nodebox/.style={draw=fg, thick, rounded corners=2pt, inner sep=5pt, text=fg},
  accentbox/.style={draw=accent, line width=1pt, rounded corners=2pt, inner sep=5pt, text=fg},
  gridcell/.style={draw=muted, minimum size=1.05cm, anchor=center, text=fg},
  treenode/.style={circle, draw=fg, thick, minimum size=0.9cm, inner sep=1pt, text=fg},
  good fill/.style={fill=good!14!white, draw=good, text=fg},
  bad fill/.style={fill=bad!12!white, draw=bad, text=fg},
  acc fill/.style={fill=accent!13!white, draw=accent, text=fg},
  % ADDED: the dashed slot a missing plate leaves behind.
  plateslot/.style={draw=muted, dashed, line width=0.8pt, rounded corners=4pt,
                    align=center, text=muted, font=\scriptsize, inner sep=6pt},
}

% ---------------------------------------------------------------------------
% ADDED for this deck: framed plate, height constrained.
% Lifted from mandatory-weekends-opening/slides/beamer-preamble.tex, where it is
% called \framedimgh and used for the Kant and Friedrich plates.
% ---------------------------------------------------------------------------
\newcommand{\framedimgh}[2]{%
  \tikz{\node[inner sep=0pt, draw=black!18, line width=0.8pt, rounded corners=4pt]%
        {\includegraphics[height=#1]{#2}};}}

% ---------------------------------------------------------------------------
% ADDED for this deck: \plateh{<height>}{<placeholder width>}{<path>}
%
% If the file is present, this is exactly \framedimgh. If it is ABSENT, it draws
% a dashed box of the same height naming the file that belongs there, so the deck
% compiles and projects with an empty img/ directory and the gap is visible
% rather than silent.
%
% This exists because no image is downloaded into this repository: the licences
% are recorded in img/IMAGE-SOURCES.txt and the files are fetched by hand, once,
% by whoever is going to project the deck. Do not "fix" this by committing a
% stand-in image; a placeholder that looks like a picture is worse than a box
% that says a picture is missing.
% ---------------------------------------------------------------------------
\newcommand{\plateh}[3]{%
  \IfFileExists{#3}{\framedimgh{#1}{#3}}{%
    \tikz{\node[plateslot, minimum height=#1, minimum width=#2]%
          {Plate not present\\[2pt]\texttt{#3}\\[2pt] See \texttt{img/IMAGE-SOURCES.txt}};}}}

% ---------------------------------------------------------------------------
% ADDED for this deck: \origslot{<original-language text>}{<work and locus>}
%
% House grammar rule 1 puts the ORIGINAL LANGUAGE on top and the English gloss
% underneath. House rule "do not put an unpinned quotation on a slide" says the
% original may only be projected once somebody has pinned it character for
% character against a primary text.
%
% Those two rules collide on this deck, because its one German-original
% quotation could NOT be pinned in the build session: two different German
% wordings circulate and neither was traceable to a printed page of
% Logik der Forschung. So the original-language slot takes the pinned text when
% there is one, and otherwise draws a dashed marker naming the work and the
% locus that still has to be checked.
%
% Pass a non-empty first argument and it renders as an ordinary quotation.
% Pass an empty first argument and the gap is visible on the slide.
%
% Do NOT "fix" an empty slot by pasting in a wording found in a quotation site.
% That is the exact failure the marker exists to prevent.
% ---------------------------------------------------------------------------
\newcommand{\origslot}[2]{%
  \if\relax\detokenize{#1}\relax
    \tikz{\node[plateslot, align=left, text width=0.86\textwidth]%
          {Original not yet pinned: #2.\\
           Pin it against a printed page before projecting, then pass it to
           \texttt{\string\origslot}. See \texttt{../SOURCES.md}.};}%
  \else
    #1%
  \fi}

% --- No closing-punchline helper here, on purpose: a frame ends on its final
%     reveal. Do not add one; a macro in the preamble reads as an instruction.
%
%     DELIBERATELY NOT DEFINED: \takeaway. The Weekend 1 FDD preamble that the
%     mandatory-weekends opening was copied from defines a centred amber-ruled
%     pull quote, and it predates the rule. Carlos, 2026-08-11: the big centred
%     blue punchline is out. A frame ends on its final reveal, and a closing
%     thought that is real content is the last ordinary line, in body text.
%     See .claude/skills/lecture-slides/SKILL.md, "No closing punchline".
%     The two closing thoughts in this deck (the placement line on frame 1 and
%     the "this hour is about the fourth" line on frame 4) are \footnotesize
%     body text at the foot of their frame, in ink, never in accent. ---


\title{Guess, Check, Correct}
\subtitle{Cold open, weekend 1, Szymborska variant}
\author{Carlos Cotrini}
\date{Friday 4 September 2026}
% No \maketitle on purpose: this deck is pasted in front of the intro and
% onboarding block, which already carries the course title page. Same decision
% as the sibling Descartes deck, and for the same reason.

\begin{document}

% ================= 1. Szymborska, and the sentence =================
% Five overlay pages. Same column grammar as the Popper frame that follows it:
% plate on the left, the placement in two short spoken lines, then the original
% language, then the gloss underneath in muted small type.
%
% WHY SHE OPENS THE DECK, ahead of the three witnesses.
% Popper, Turing and Whewell are all making a claim ABOUT method, from outside.
% Szymborska makes no claim at all: she describes what it is like from the
% inside to finish something, read it back, and find that it will not hold.
% Every person in that room at 08:00 has done that. Nobody in it has falsified
% a conjecture. So she is the frame that buys the other three their audience.
%
% THE HAND-OFF IS THE WHOLE POINT, AND IT IS FRAGILE. This frame ends on
%     rodzi sie z bezustannego "nie wiem"
%     born from a continuous "I don't know"
% and the very next frame opens on
%     We do not know: we can only guess.
% The room hears the same sentence twice inside twenty seconds, once from a poet
% and once from a philosopher, and the second one adds the guess. DO NOT put a
% frame between these two and do not reorder them. If the morning runs late,
% cut overlay 4 (the repeat), never the last overlay.
%
% ORIGINAL LANGUAGE, and this frame is better off than its neighbours.
% The lecture was delivered in Polish at the Swedish Academy on 7 December 1996
% and nobelprize.org carries the Polish and the English side by side. So unlike
% the Popper frame, this frame's \origslot takes REAL PINNED TEXT and draws no
% dashed marker: the Polish was read off the Polish page in the build session
% and the English off the English page. Both are in ../SOURCES.md.
%
% A TRANSLATION NOTE THAT MATTERS. The Polish is SINGULAR, "Poeta rowniez,
% jesli jest prawdziwym poeta" (the poet too, if he is a true poet), and
% Baranczak and Cavanagh pluralise it to "Poets, if they're genuine". The gloss
% on this slide is their published wording and NOT a correction of it. Anyone
% comparing the two lines will see the number change; that is the translators'
% choice, and ../SOURCES.md says so.
%
% ABRIDGEMENT. Each quotation is cut once and the cut is marked \ldots.
% The full pinned text, for the record and for anyone restoring it:
%
%  PL  Poeta rowniez, jesli jest prawdziwym poeta, musi ciagle powtarzac sobie
%      "nie wiem". Kazdym utworem probuje na to odpowiedziec, ale kiedy tylko
%      postawi kropke, juz ogarnia go wahanie, juz zaczyna sobie zdawac sprawe,
%      ze jest to odpowiedz tymczasowa i absolutnie niewystarczajaca. Wiec
%      probuje jeszcze raz i jeszcze raz, a potem te kolejne dowody jego
%      niezadowolenia z siebie historycy literatury zepna wielkim spinaczem i
%      nazywac beda "dorobkiem" ...
%
%  EN  Poets, if they're genuine, must also keep repeating "I don't know." Each
%      poem marks an effort to answer this statement, but as soon as the final
%      period hits the page, the poet begins to hesitate, starts to realize that
%      this particular answer was pure makeshift that's absolutely inadequate to
%      boot. So the poets keep on trying, and sooner or later the consecutive
%      results of their self-dissatisfaction are clipped together with a giant
%      paperclip by literary historians and called their "oeuvre" ...
%
% THE CLAUSE THIS DECK DELIBERATELY DOES NOT PROJECT is the paperclip. The poet
% keeps every past attempt and files it as an oeuvre; the algorithm at 09:00
% keeps only the current guess and throws the rest away. That is a real
% disanalogy and it is the one thing here worth saying out loud IF SOMEBODY
% ASKS. It stays off the slide, because a cold open that argues with itself is
% four frames too long.
%
% POLISH DIACRITICS are written as LaTeX accent commands and not as raw UTF-8,
% which is the convention the sibling Descartes deck already uses for its
% French. l with stroke is \l{}, the ogonek is \k{}, the dot above is \.{}.
% That keeps this file byte-clean under pdflatex with T1 and no assumption
% about inputenc.
%
% QUOTATION MARKS. nobelprize.org sets the Polish with English-style curly
% quotes rather than the Polish low-opening form, and this frame reproduces
% what was actually pinned. Do not "correct" it to ,,nie wiem''.
%
% THE ALERTS ARE THE SAME CODE AS THE REST OF THE DECK: exactly the clause that
% carries the loop, and nothing else. Doubt on overlay 3, repetition on overlay
% 4, the not-knowing on overlay 5. Three alerts, three beats, one per click.
\begin{frame}{Szymborska, 1996: the answer that will not hold}
  \vspace{0mm}
  \begin{columns}[c]
    \column{0.30\textwidth}\centering
      \plateh{3.6cm}{3.0cm}{img/szymborska-kubik.jpg}\\[3pt]
      % The break is FORCED. Left to itself this 0.30 column split the licence
      % as "CC / BY 3.0", which the 130 dpi render caught and the log did not.
      {\scriptsize\color{muted}Wis\l{}awa Szymborska, Krak\'ow, 2009.\\[1pt]
       Mariusz Kubik, CC BY 3.0}
    \column{0.66\textwidth}
      \onslide<2->{{\footnotesize Stockholm, 7 December 1996. Szymborska is
      asked to say what a poet does.}}\\[6pt]
      % NB: the break between the Polish and its gloss is "\\[2pt]", NEVER
      % "\par". A \par inside these groups ends the paragraph, and the "\\[8pt]"
      % that follows the group then fails with "There's no line here to end".
      % That is what this frame did on its first build.
      \onslide<3->{{\footnotesize Natchnienie, czymkolwiek ono jest, rodzi
       si\k{e} z \alert{bezustannego ``nie wiem''}.\\[2pt]
       {\scriptsize\color{muted}Whatever inspiration is, it's born from a
        continuous ``I don't know.''}}}\\[6pt]
      \onslide<4->{{\footnotesize Ka\.zdym utworem pr\'obuje na to
       odpowiedzie\'c, ale kiedy tylko \alert{postawi kropk\k{e}, ju\.z ogarnia
       go wahanie} \ldots\ Wi\k{e}c pr\'obuje \alert{jeszcze raz i jeszcze
       raz}.\\[2pt]
       {\scriptsize\color{muted}Each poem marks an effort to answer this
        statement, but as soon as the final period hits the page, the poet
        begins to hesitate \ldots\ So the poets keep on trying.}}}\\[4pt]
      {\scriptsize\color{muted}Wis\l{}awa Szymborska, ``The poet and the
       world'', Nobel Lecture 1996. Polish original; English by Bara\'nczak and
       Cavanagh. Copyright The Nobel Foundation 1996.}
  \end{columns}
\end{frame}

% ================= 2. Popper, and the sentence =================
% Four overlay pages. Left column is the plate. Right column lands the placement
% in two short spoken lines, then the original-language slot, then the sentence.
%
% The frame ends on the sentence. There is no closing line under it.
%
% THIS IS THE FRAME THAT CARRIES THE ROOM. Everything after it is evidence.
\begin{frame}{Popper, 1934: conjecture and refutation}
  \vspace{1mm}
  \begin{columns}[c]
    \column{0.32\textwidth}\centering
      \plateh{4.0cm}{3.0cm}{img/popper-lse.jpg}\\[3pt]
      {\scriptsize\color{muted}Karl Popper, about 1980.\\LSE Library}
    \column{0.64\textwidth}
      \onslide<2->{{\footnotesize Vienna, 1934. Popper asks what entitles anyone
      to believe a general rule drawn from finitely many cases.}}\\[6pt]
      \onslide<2->{{\footnotesize His answer is that nothing does, and that this
      is survivable.}}\\[10pt]
      \onslide<3->{\origslot{}{\textit{Logik der Forschung} (1934), Abschnitt 85}}\\[6pt]
      \onslide<4->{{\large We do not know: \alert{we can only guess.}}}\\[8pt]
      {\scriptsize\color{muted}Karl Popper, \textit{The Logic of Scientific
       Discovery}, chapter 10, section 85 (1959), the authorised English of
       % Ordinary citation. It carried "Located, not yet pinned: a colon and
       % a semicolon version both circulate. See ../SOURCES.md." until
       % 2026-09-02, when Carlos struck that whole class of note from every
       % deck: nothing addressed to the builder goes on a slide. The caveat
       % lives in ../SOURCES.md and in the comment above this frame.
       \textit{Logik der Forschung} (1934).}
  \end{columns}
\end{frame}

% ================= 3. The tetradic schema =================
% Seven overlay pages. The pinned Conjectures and Refutations sentence at the
% top, then the four nodes one per click, then the return arrow, then one wave
% that hangs Friday morning off each node.
%
% ABRIDGEMENT. The quotation is cut in two places and each cut is marked \ldots.
% The full pinned sentence, for the record:
%
%   The way in which knowledge progresses, and especially our scientific
%   knowledge, is by unjustified (and unjustifiable) anticipations, by guesses,
%   by tentative solutions to our problems, by conjectures. These conjectures
%   are controlled by criticism; that is, by attempted refutations, which
%   include severely critical tests.
%
% THE DIAGRAM IS DRAWN WITH "to", NEVER "--". A TikZ line segment is the one
% place the house dash rule allows a double hyphen, and this deck declines even
% that, so that the dash grep over this source returns nothing at all and the
% next person does not have to eyeball hits.
%
% NO-SHIFT RULE. Every node and every arrow uses "visible on=<n->" rather than
% \only, so the picture's bounding box is identical on all seven overlays and
% nothing moves as pieces land. SKILL.md Pattern F.
%
% THE RETURN ARROW GOES OVER THE TOP, not underneath. Underneath is where the
% two caption lines live, and an arc routed below collided with them.
\begin{frame}{Popper's tetradic schema}
  \vspace{-3mm}
  {\footnotesize The way in which knowledge progresses \ldots\ is by
   \alert{unjustified (and unjustifiable) anticipations, by guesses}, by
   tentative solutions to our problems, by conjectures. These conjectures are
   controlled by criticism; that is, \alert{by attempted refutations} \ldots}

  \vspace{1mm}
  \centering
  % hyphenpenalty is set OUTSIDE the picture so the caption nodes inherit it:
  % without it "Square it. Compare with 2026" broke as "Com-pare" inside a
  % 2.9cm node. Same trick the sibling deck uses on its precept tabular.
  % ... and it is SCOPED to the picture: left switched on to the end of the
  % frame it also suppressed hyphenation in the source line below, which then
  % ran 4.45pt past the right margin.
  {\hyphenpenalty=10000 \exhyphenpenalty=10000
  \begin{tikzpicture}
    % --- the chain ---
    \node[nodebox, visible on=<2->] (p1) at (0,0)    {$P_1$};
    \node[nodebox, visible on=<3->] (tt) at (3.5,0)  {TT};
    \node[nodebox, visible on=<4->] (ee) at (7.0,0)  {EE};
    \node[nodebox, visible on=<5->] (p2) at (10.5,0) {$P_2$};
    \draw[edge, visible on=<3->] (p1) to (tt);
    \draw[edge, visible on=<4->] (tt) to (ee);
    \draw[edge, visible on=<5->] (ee) to (p2);
    % --- what each letter stands for ---
    \node[below=1.5mm of p1, visible on=<2->] (n1) {Problem};
    \node[below=1.5mm of tt, visible on=<3->] (n2) {Tentative theory};
    \node[below=1.5mm of ee, visible on=<4->] (n3) {Error elimination};
    \node[below=1.5mm of p2, visible on=<5->] (n4) {New problem};
    % --- the loop closes ---
    \draw[dimedge, visible on=<6->] (p2.north)
      to[out=90, in=90, looseness=0.26]
      node[midway, above, font=\scriptsize, text=muted] {and again} (p1.north);
    % --- Friday morning, hung off the schema, as ONE wave ---
    \node[below=1mm of n1, text width=3.0cm, align=center, text=muted,
          font=\scriptsize, visible on=<7->] {You cannot write the answer down};
    \node[below=1mm of n2, text width=3.0cm, align=center, text=muted,
          font=\scriptsize, visible on=<7->] {A first guess, $\theta = 45$};
    \node[below=1mm of n3, text width=3.0cm, align=center, text=muted,
          font=\scriptsize, visible on=<7->] {Square it, compare with 2026};
    \node[below=1mm of n4, text width=3.0cm, align=center, text=muted,
          font=\scriptsize, visible on=<7->] {A smaller miss, and a better question};
  \end{tikzpicture}}

  \vspace{0mm}
  {\raggedright\scriptsize\color{muted}\hspace*{0.2cm}Schema after Karl Popper,
   \textit{Objective Knowledge}, chapter 8 (1972). Quotation: \textit{Conjectures
   and Refutations}, preface (1963).\par}
\end{frame}

% ================= 4. Turing =================
% Four overlay pages. Two quotations, both pinned, then the paraphrase.
%
% The second quotation is the bridge to the 11:00 lecture: a search over
% behaviours, scored by a criterion, is gradient descent with the words changed.
% If the morning runs late, cut overlay 3 rather than the frame.
\begin{frame}{Turing, 1950: the child machine}
  \vspace{1mm}
  \begin{columns}[c]
    \column{0.32\textwidth}\centering
      \plateh{4.0cm}{3.0cm}{img/turing-1951.jpg}\\[3pt]
      {\scriptsize\color{muted}Alan Turing, 29 March 1951.\\Elliott and Fry}
    \column{0.64\textwidth}
      \onslide<2->{{\footnotesize We cannot expect to find a good child-machine
      at the first attempt. One must experiment with teaching one such machine
      and see how well it learns. One can then \alert{try another} and see if it
      is better or worse.}}\\[10pt]
      \onslide<3->{{\footnotesize Now the learning process may be regarded as
      \alert{a search for a form of behaviour} which will satisfy the teacher
      (or some other criterion).}}\\[10pt]
      \onslide<4->{{\scriptsize\color{muted}You do not get it right the first
      time. You try, you measure, you try again.}}\\[8pt]
      {\scriptsize\color{muted}A. M. Turing, ``Computing Machinery and
       Intelligence'', \textit{Mind} LIX(236), 1950, section 7, Learning
       Machines.}
  \end{columns}
\end{frame}

% ================= 5. Whewell =================
% Four overlay pages. Same shape as frame 3.
%
% THE SECOND QUOTATION IS THE ONE THAT PAYS FOR THE FRAME. "Overshooting the
% mark" is what a step size does at 11:00, written in 1858 about inventive minds
% in general. Do not cut it in favour of the first.
%
% MISQUOTATION TRAP, found in the research session and worth keeping in mind
% every time this line is retyped: the version in general circulation reads
% "To try wrong guesses is apparently the only way to hit upon right ones."
% The 1858 text has no "apparently". It reads "is, with most persons,".
% The wording below was read off the third-edition text, not off a quotation
% site. Do not "fix" it back.
%
% ALTERNATIVE QUOTATION, if Carlos prefers the validation reading to the
% guessing one, from the aphorisms at the head of the same chapter:
%   Aphorism XII. It is a test of true theories not only to account for, but to
%   predict phenomena.
% That one reaches 16:00 and Saturday's validation block rather than the
% morning. It is NOT used here because this deck's job is the morning loop.
\begin{frame}{Whewell, 1858: trying wrong guesses}
  \vspace{1mm}
  \begin{columns}[c]
    \column{0.32\textwidth}\centering
      \plateh{4.0cm}{3.0cm}{img/whewell-wellcome.jpg}\\[3pt]
      {\scriptsize\color{muted}William Whewell, stipple engraving.\\Wellcome
       Collection, CC BY 4.0}
    \column{0.64\textwidth}
      \onslide<2->{{\footnotesize To try \alert{wrong guesses} is, with most
      persons, the only way to hit upon right ones.}}\\[10pt]
      \onslide<3->{{\footnotesize The tendencies of our speculative nature
      \ldots\ perpetually show their vigour by \alert{overshooting the mark}.
      They obtain something, by aiming at much more.}}\\[10pt]
      \onslide<4->{{\scriptsize\color{muted}Most guesses are wrong, and a step
      that aims too far is how the right one gets found.}}\\[8pt]
      {\scriptsize\color{muted}William Whewell, \textit{Novum Organon
       Renovatum}, third edition (1858), book II, chapter V, page 79.}
  \end{columns}
\end{frame}

% ================= 6. References =================
% ONE overlay page, static, and deliberately so: this is a handout page and a
% rights record rather than teaching content, and SKILL.md hard rule 5 is a
% teaching-content rule. It is NOT part of the three minutes.
%
% Every image line carries its licence because two of the three are not plain
% public domain: the Wellcome scan is CC BY 4.0 and requires the credit that is
% printed on frame 4, and the Popper photograph carries a Flickr Commons
% "no known copyright restrictions" statement, which is an institutional
% assertion rather than a licence. See img/IMAGE-SOURCES.txt.
\begin{frame}{References}
  \vspace{-3mm}
  {\scriptsize
  \begin{columns}[T]
    \column{0.53\textwidth}
      Texts\\[4pt]
      {\color{muted}
      Szymborska, Wis\l{}awa. ``The poet and the world''. Nobel Lecture,
      Stockholm, 7 December 1996. Polish and English (tr. Bara\'nczak and
      Cavanagh). Copyright The Nobel Foundation 1996.\\[3pt]
      Popper, Karl R. \textit{The Logic of Scientific Discovery}. London:
      Hutchinson, 1959. English of \textit{Logik der Forschung}, Vienna:
      Springer, 1934. Chapter 10, section 85.\\[3pt]
      Popper, Karl R. \textit{Conjectures and Refutations: The Growth of
      Scientific Knowledge}. London: Routledge and Kegan Paul, 1963. Preface.\\[3pt]
      Popper, Karl R. \textit{Objective Knowledge: An Evolutionary Approach}.
      Oxford: Clarendon Press, 1972. Chapter 8.\\[3pt]
      Turing, A. M. ``Computing Machinery and Intelligence''. \textit{Mind}
      LIX(236), 1950, pages 433 to 460. Section 7.\\[3pt]
      Whewell, William. \textit{Novum Organon Renovatum}. Third edition of
      \textit{The Philosophy of the Inductive Sciences}, part two. London:
      John W. Parker, 1858. Book II, chapter V.}
    \column{0.43\textwidth}
      Images\\[4pt]
      {\color{muted}
      Wis\l{}awa Szymborska, Krak\'ow, 2009. Mariusz Kubik. CC BY 3.0 and
      GFDL 1.2 or later.\\[3pt]
      Karl Popper, about 1980. LSE Library, via Flickr Commons. Marked
      ``no known copyright restrictions''.\\[3pt]
      Alan Turing, 29 March 1951. Elliott and Fry. Public domain.\\[3pt]
      William Whewell, stipple engraving. Wellcome Collection L0014766,
      CC BY 4.0.}\\[5pt]
      Rights note\\[3pt]
      {\color{muted}Popper died in 1994 and Szymborska in 2012, so their texts
      are protected in Switzerland until the end of 2064 and 2082. The
      quotations run on the quotation right, Art. 25 URG: a published work, an
      extent justified by the purpose, the quotation marked and its source
      named.}
  \end{columns}}
\end{frame}

\end{document}
